\par\Needspace{10\baselineskip}\pdfbookmark[1]{Table for more convenient calculation. 1812.}{gauss1812log}
\begin{center}
\rule{0.66\linewidth}{0.3pt}\\
Monatliche Correspondenz, edited by Baron von \textsc{Zach}. 1812. Nov.\\
\rule{0.66\linewidth}{0.3pt}
\end{center}

\emph{Table for more conveniently calculating the logarithm of the sum or difference of two quantities that are themselves given only by their logarithms. By Professor} \textsc{Gauss}.

As the work of astronomers engaged in computation continues to expand, every simplification, however small in itself, becomes more important to them. The \emph{Monatliche Correspondenz} has already rendered many services in this respect by publishing various tables too small to be issued separately. I am therefore glad to contribute to it a small table which, though not strictly astronomical, will nevertheless be especially welcome to astronomers engaged in computation, and which might in future very suitably accompany a new printing of \textsc{La
Lande}'s small tables. The operation it is intended to simplify occurs at every turn in astronomical calculations. It otherwise requires three lookups in the logarithm tables, or, with an easy transformation, still necessarily two; here these are reduced to one. As far as I know, \textsc{Leonelli} was the first to propose the idea. But he intended to construct such a table for calculations with 14 decimal places, and it is precisely this that I cannot consider appropriate. On that scale it would fill a large folio volume, calculating it would demand an enormous amount of work and time, and it would hardly ever be useful, and always only of little use. Calculations of such accuracy occur so rarely---in practical astronomy proper, never---that the relatively small saving would not repay the construction, or even the purchase, of such a table. Many years ago I constructed this table for my own use in calculations with 5 decimal places, which are the most frequent in practice. The saving, though small \emph{each time}, becomes very considerable when repeated many thousands of times, and has already amply repaid the effort I put into it. It would be desirable for someone to undertake the work of constructing a similar table on 10 or 100 times the scale for calculations with 7 decimal places, \SourceLeaf{254}
which could be added to the ordinary logarithm tables as a very valuable supplement.

\Needspace{24\baselineskip}The arrangement of the table, which consists of three columns, is very simple. The first column \(A\) runs from 0 to 2 by every thousandth, from there to \(3,4\) by every hundredth, and from \(3,4\) to \(5,0\) by every tenth. With \(5,0\) the table can be considered finished for 5 decimal places, since the second column vanishes for this and for larger values of \(A\), and the numbers in the third column become equal to those in the first. If a number in the first column is written \(A=\log m\), the second column gives \(B=\log(1+\frac{1}{m})\) and the third \(C=\log(1+m)\), so that always \(C=A+B\). The numbers in the three columns can therefore also be regarded as twice the logarithms of the tangents, cosecants and secants of angles from \(45^{0}\) to \(90^{0}\). Their use is as follows:

\begin{center}
I.\quad \emph{To find, from the logarithms of two quantities \(a,b\), the logarithm of their sum.}
\end{center}

Let \(\log a\) be the larger logarithm. One enters \(\log a-\log b\) in column \(A\) and takes alongside it either \(B\) from the second column or \(C\) from the third. Then
\[
\begin{aligned}
&\log(a+b)=\log a+B\\
\text{or}\qquad &\log(a+b)=\log b+C.
\end{aligned}
\]

\begin{center}
II.\quad \emph{To find, from the logarithms of two quantities \(a,b\), the logarithm of their difference.}
\end{center}

\emph{First,} if the difference of the logarithms \(\log a-\log b\) is greater than \(0,30103\), it is looked up in \(C\). This gives
\[
\begin{aligned}
&\log(a-b)=\log a-B\\
\text{or}\qquad &\log(a-b)=\log b+A
\end{aligned}
\]

\emph{Second,} if \(\log a-\log b\) is smaller than \(0,30103\), it is looked up in \(B\). This gives
\[
\begin{aligned}
&\log(a-b)=\log a-C\\
\text{or}\qquad &\log(a-b)=\log b-A
\end{aligned}
\]
\SourceLeaf{255}
There are therefore two ways of solving each problem. It is advisable, however, to become accustomed to one particular way, so that using the table can more easily become automatic. I have found the method given first in each case the most convenient for this.

\Needspace{12\baselineskip}\begin{center}\textit{Examples:}\end{center}

I. To find the logarithm of the sum from $\log a=0{,}36173$ and $\log b=0{,}23045$, $0{,}13128$ is looked up in $A$. This gives
\[
\begin{array}{rcr}
B&\cdot\;\cdot\;\cdot\;\cdot\;\cdot&0{,}24033\\
\log a&\cdot\;\cdot\;\cdot\;\cdot\;\cdot&0{,}36173\\\hline
\log(a+b)&\cdot\;\cdot&0{,}60206
\end{array}
\qquad
\begin{array}{rcr}
C&\cdot\;\cdot\;\cdot\;\cdot\;\cdot&0{,}37161\\
\log b&\cdot\;\cdot\;\cdot\;\cdot\;\cdot&0{,}23045\\\hline
&&0{,}60206
\end{array}
\]

II. The logarithm of the difference is required from $\log a=0{,}89042$ and $\log b=0{,}24797$. Since $\log a-\log b=0{,}64245$ is greater than $0{,}30103$, it is looked up in column $C$, alongside which are found
\[
\begin{array}{rcr}
B&\cdot\;\cdot\;\cdot\;\cdot\;\cdot&0{,}11227\\
\log a&\cdot\;\cdot\;\cdot\;\cdot\;\cdot&0{,}89042\\\hline
\log(a-b)&\cdot\;\cdot&0{,}77815
\end{array}
\qquad
\begin{array}{rcr}
A&\cdot\;\cdot\;\cdot\;\cdot\;\cdot&0{,}53018\\
\log b&\cdot\;\cdot\;\cdot\;\cdot\;\cdot&0{,}24797\\\hline
&&0{,}77815
\end{array}
\]

III. The logarithm of the difference is required from $\log a=0{,}25042$ and $\log b=0{,}19033$. Here $\log a-\log b=0{,}06009$, looked up in $B$, gives
\[
\begin{array}{rcr}
C&\cdot\;\cdot\;\cdot\;\cdot\;\cdot&0{,}88871\\
\log a&\cdot\;\cdot\;\cdot\;\cdot\;\cdot&0{,}25042\\\hline
\log(a-b)&\cdot\;\cdot&9{,}36171
\end{array}
\qquad
\begin{array}{rcr}
A&\cdot\;\cdot\;\cdot\;\cdot\;\cdot&0{,}82862\\
\log b&\cdot\;\cdot\;\cdot\;\cdot\;\cdot&0{,}19033\\\hline
&&9{,}36171
\end{array}
\]

